\documentclass[10pt,a4paper]{amsart}
\usepackage{amsmath,amssymb,amsthm,graphicx}
\usepackage{verbatim}
\newenvironment{reference}{\comment}{\endcomment}
%\newenvironment{reference}{}{}
\newenvironment{slogan}{\comment}{\endcomment}
\newenvironment{history}{\comment}{\endcomment}

% For commutative diagrams we use Xy-pic
\usepackage[all]{xy}

% We use 2cell for 2-commutative diagrams.
\xyoption{2cell}
% Theorem environments.
%
\theoremstyle{plain}
\newtheorem{theorem}[subsection]{Theorem}
\newtheorem{proposition}[subsection]{Proposition}
\newtheorem{lemma}[subsection]{Lemma}

\theoremstyle{definition}
\newtheorem{definition}[subsection]{Definition}
\newtheorem{example}[subsection]{Example}
\newtheorem{exercise}[subsection]{Exercise}
\newtheorem{situation}[subsection]{Situation}

\theoremstyle{remark}
\newtheorem{remark}[subsection]{Remark}
\newtheorem{remarks}[subsection]{Remarks}

\numberwithin{equation}{subsection}

% Macros
%
\def\lim{\mathop{\mathrm{lim}}\nolimits}
\def\colim{\mathop{\mathrm{colim}}\nolimits}
\def\Spec{\mathop{\mathrm{Spec}}}
\def\Hom{\mathop{\mathrm{Hom}}\nolimits}
\def\Ext{\mathop{\mathrm{Ext}}\nolimits}
\def\SheafHom{\mathop{\mathcal{H}\!\mathit{om}}\nolimits}
\def\SheafExt{\mathop{\mathcal{E}\!\mathit{xt}}\nolimits}
\def\Sch{\mathit{Sch}}
\def\Mor{\mathop{\mathrm{Mor}}\nolimits}
\def\Ob{\mathop{\mathrm{Ob}}\nolimits}
\def\Sh{\mathop{\mathit{Sh}}\nolimits}
\def\NL{\mathop{N\!L}\nolimits}
\def\CH{\mathop{\mathrm{CH}}\nolimits}
\def\proetale{{pro\text{-}\acute{e}tale}}
\def\etale{{\acute{e}tale}}
\def\QCoh{\mathit{QCoh}}
\def\Ker{\mathop{\mathrm{Ker}}}
\def\Im{\mathop{\mathrm{Im}}}
\def\Coker{\mathop{\mathrm{Coker}}}
\def\Coim{\mathop{\mathrm{Coim}}}

% Boxtimes
%
\DeclareMathSymbol{\boxtimes}{\mathbin}{AMSa}{"02}

%
% Macros for moduli stacks/spaces
%
\def\QCohstack{\mathcal{QC}\!\mathit{oh}}
\def\Cohstack{\mathcal{C}\!\mathit{oh}}
\def\Spacesstack{\mathcal{S}\!\mathit{paces}}
\def\Quotfunctor{\mathrm{Quot}}
\def\Hilbfunctor{\mathrm{Hilb}}
\def\Curvesstack{\mathcal{C}\!\mathit{urves}}
\def\Polarizedstack{\mathcal{P}\!\mathit{olarized}}
\def\Complexesstack{\mathcal{C}\!\mathit{omplexes}}
% \Pic is the operator that assigns to X its picard group, usage \Pic(X)
% \Picardstack_{X/B} denotes the Picard stack of X over B
% \Picardfunctor_{X/B} denotes the Picard functor of X over B
\def\Pic{\mathop{\mathrm{Pic}}\nolimits}
\def\Picardstack{\mathcal{P}\!\mathit{ic}}
\def\Picardfunctor{\mathrm{Pic}}
\def\Deformationcategory{\mathcal{D}\!\mathit{ef}}
\usepackage{cleveref}
\crefname{originalreference}{original reference}{original references}
\Crefname{originalreference}{Original reference}{Original references}
\makeatletter
\newlabel{lemma-degree-minimal-polynomial}{{09GN}{source}{}{sourcepointer-1}{}}
\newlabel{lemma-degree-minimal-polynomial@cref}{{[originalreference][0][]09GN}{[1][0][]source}{}{}{}}
\makeatother
\crefname{originalreference}{original reference}{original references}
\Crefname{originalreference}{Original reference}{Original references}
\makeatletter
\newlabel{lemma-multiplicativity-degrees}{{09G9}{source}{}{sourcepointer-2}{}}
\newlabel{lemma-multiplicativity-degrees@cref}{{[originalreference][0][]09G9}{[1][0][]source}{}{}{}}
\makeatother
\title{The characteristic polynomial of a field element as a linear map is a power of its minimal polynomial}
\author{The Stacks Project Authors}
\date{stacks-project:fields.tex:local-snapshot-20261012; GNU FDL 1.2 or later}
\begin{document}
\maketitle
\noindent\textit{Editorial note.} All delivered mathematical text is the author's own, copied from the pinned Stacks Project chapter file fields.tex; no statement or proof step was rewritten, reordered or removed.\par\medskip
\noindent\textit{Editorial note.} The counted proof is the author's complete argument: This lemma feeds lemma-trace-and-norm-from-minimal-polynomial (source lines 2642-2658) and the trace-pairing separability criterion (delivered L48333).\par\medskip
\noindent\textit{Editorial note.} The label lemma-characteristic-vs-minimal-polynomial names the counted statement, and the proof environment immediately following it in the source is delivered verbatim as its proof.\par\medskip
\noindent\textit{Editorial note.} The counted statement is carried by the author's own lemma environment in the chapter source.\par\medskip
\noindent\textit{Allowed standard prerequisites (editorial).}\begin{itemize}
\item Same chapter: lemma-degree-minimal-polynomial (deg of the minimal polynomial equals the subextension degree) and lemma-multiplicativity-degrees (delivered card L48338) for e deg(P) = [L : K]; Cayley-Hamilton; minimal polynomial divides the characteristic polynomial
\end{itemize}\medskip
\ifdefined\hypertarget\hypertarget{sourcepointer-1}{}\fi
\noindent\textit{Source reference.} Original reference 09GN: sources/stacks-fields/pinned-original-source.tex, source lines 1087--1087. Prerequisite pointer invoked at 2626; it names a standard fact of the same chapter that the counted argument uses, and every step of the counted proof built on it is retained verbatim in the delivered body.\par\medskip
\ifdefined\hypertarget\hypertarget{sourcepointer-2}{}\fi
\noindent\textit{Source reference.} Original reference 09G9: sources/stacks-fields/pinned-original-source.tex, source lines 759--759. Prerequisite pointer invoked at 2627; it names a standard fact of the same chapter that the counted argument uses, and every step of the counted proof built on it is retained verbatim in the delivered body.\par\medskip
\begin{lemma}
\label{lemma-characteristic-vs-minimal-polynomial}
Let $L/K$ be a finite extension of fields. Let $\alpha \in L$ and let $P$
be the minimal polynomial of $\alpha$ over $K$. Then the characteristic
polynomial of the $K$-linear map $\alpha : L \to L$ is equal to
$P^e$ with $e \deg(P) = [L : K]$.
\end{lemma}

\begin{proof}
Choose a basis $\beta_1, \ldots, \beta_e$ of $L$ over $K(\alpha)$.
Then $e$ satisfies $e \deg(P) = [L : K]$ by
Lemmas \ref{lemma-degree-minimal-polynomial} and
\ref{lemma-multiplicativity-degrees}.
Then we see that $L = \bigoplus K(\alpha) \beta_i$ is a
direct sum decomposition into $\alpha$-invariant subspaces
hence the characteristic polynomial of $\alpha : L \to L$
is equal to the characteristic polynomial of
$\alpha : K(\alpha) \to K(\alpha)$ to the power $e$.

\medskip\noindent
To finish the proof we may assume that $L = K(\alpha)$.
In this case by Cayley-Hamilton we see that $\alpha$
is a root of the characteristic polynomial. And since the
characteristic polynomial has the same degree as the minimal
polynomial, we find that equality holds.
\end{proof}

\end{document}
